\scanpage{27}
% L1–3
$\unsure{H^1|H\!H}(X,B_X)$ le sous-espace vectoriel de l'espace de dim finie $H^1(X,\underline O_X)$, \deleted{\unsure{an|annulé}} noyau de Frobenius.

% L4–7
\uline{Remarques} Soit $A=k[[X]]$, où $k$ est un corps algébriquement clos. Les revêtements \uline{non} triviaux \margin{\unsure{les|ses}} \deleted{ramifiés} de $A$ obtenus par extraction d'une racine $p$ième de $A$ sont \unsure{essentiellement|essentielt} isomorphes
% L8–9
[à condition de changer l'uniformisante, au besoin] à
\[
A[X^{n/p}]\simeq A[T]/(T^p-X^n)
\]
% L10–13
où $\unsure{1|0}\leq n\leq p-1$ \deleted{\unsure{si}} [pour $n=0$, on obtiendrait $\simeq A[T]/T^pA[\unsure{t|T}]$]. Si $n=1$, c'est un anneau local régulier. Sinon, il ne l'est pas, comme on voit facilement.

% L14–16
Par exemple, si $p=3$, \deleted{\unsure{si $A$}} et si on prend $n=\unsure{2|1}$, on trouve
\[
A[T]/(T^{\unsure{3|2}}-X^2)\simeq A[X^{2/3}],
\]
% L17–21
qui est isomorphe au sous-anneau de $\deleted{\unsure{B=}}k[[Y]]=k[[X^{1/3}]]$ formé des séries formelles dont le coefficient de $Y$ est nul. Pourtant, du point de vue de la théorie des $B$-\deleted{fibrés}\inserted{revêtements},
% archive sheet number: 26
